\section{Five Claims That Failed}
\label{sec:negative}

We designed this model around four mechanisms and wrote down what each should
produce before testing it. All four failed, and so did a fifth claim we
expected to survive. We report them with the tests that produced them, because
each is informative about the problem rather than only about our design, and
because a reviewer with our artifact would find them in an afternoon.

\subsection{Few-shot enrollment: the curve is flat}

The deployed embedding is trained with the same metric objective used at
inference, which should let a new application be enrolled from a handful of
labelled flows by averaging them into a prototype. We pre-registered the test:
\emph{if enrollment at $k = 5$--$25$ does not close the gap to the best
baseline, the architecture has no case on this dataset.}

Service-level macro-F1 moves from 0.502 $\pm$ 0.070 at $k=1$ to 0.505 $\pm$
0.070 at $k=100$: $+0.003$ across two orders of magnitude in labelled data,
against a seed spread twenty times larger. The gap does not close.

The flat curve is not a mechanical failure. The prototype is \emph{already
converged at $k=1$}, which says the embedding places same-class flows tightly.
The problem is where it places them. Cheap enrollment of a weak classifier is
not a contribution.

\subsection{Temporal drift: nothing to measure}

Trained on the earliest 70\% of CESNET-QUIC22 and evaluated day by day over the
held-out tail, macro-F1 \emph{rises} from 0.768 to 0.791. There is no
degradation. The honest reading is not that the method resists drift; it is
that a four-week corpus with a six-day held-out tail cannot exhibit the
phenomenon. Reporting a flat curve as drift resistance would be read as a claim
about deployment lifetime, which this evidence cannot support. The claim is
withdrawn pending a corpus spanning months.

\subsection{Adversarial nuisance removal: it removes nothing}

We replaced an earlier invariance score --- which we found was maximised by a
constant encoder, and whose nuisance label was a threshold on two of the
model's own input features, reconstructible with agreement 1.000 --- with
standard representation probing: freeze the embedding, train linear and MLP
probes to predict the nuisance, and report accuracy above the majority-class
floor against a raw-feature control.

\begin{center}
\small
\begin{tabular}{llrr}
\toprule
Corpus & Nuisance & From $z_{\mathrm{flow}}$ & From raw features \\
\midrule
CESNET & week & $+0.0006 \pm 0.0006$ & $-0.0049 \pm 0.0041$ \\
5G & capture session & $+0.2935 \pm 0.0205$ & $+0.2790 \pm 0.0854$ \\
\bottomrule
\end{tabular}
\end{center}

On CESNET the result is vacuous: week is not decodable from the raw features
either, so there was nothing to remove, and the near-zero leak from the
embedding measures the corpus rather than the method. On 5G the nuisance is
genuinely present and the embedding leaks it \emph{at the same rate as its own
input}. The component achieves nothing.

The raw-feature control is what makes this readable. Without it, the CESNET row
alone would have supported a published claim of near-perfect invariance.

\subsection{Open-set rejection: right result, wrong mechanism}

% Generated by scripts/make_paper_assets.py. Do not edit by hand.
\begin{table}[t]
\centering
\small
\begin{tabular}{lcc}
\toprule
Rejection rule & AUROC $\uparrow$ & FPR@95TPR $\downarrow$ \\
\midrule
Prototype cosine \textbf{(ours)} & 0.919 $\pm$ 0.033 & 0.341 $\pm$ 0.107 \\
Max softmax probability & 0.837 $\pm$ 0.085 & 0.610 $\pm$ 0.136 \\
Energy score & 0.870 $\pm$ 0.081 & 0.375 $\pm$ 0.138 \\
Mahalanobis distance & 0.927 $\pm$ 0.030 & 0.316 $\pm$ 0.133 \\
\bottomrule
\end{tabular}
\caption{Open-set rejection of applications never seen in training, mean $\pm$ std over 3 seed(s). 7,509 unknown against 25,350 known test flows, held-out applications: ms\_teams, teamfight\_tactics, geforce\_now. All four rules score the \emph{same} frozen embedding, so the comparison isolates the rejection rule rather than the representation. FPR@95TPR is the column to read: the AUROC gap between prototype and softmax is within the seed spread, while the false-accept gap at matched true-positive rate is roughly twice it. Mahalanobis on the same embedding is statistically indistinguishable from the prototype rule, so the claim is that \emph{distance-based rejection beats softmax thresholding}, not that ours is the best distance rule.}
\label{tab:open-set}
\end{table}


This is the claim we expected to survive, and half of it does. Rejecting
unknown applications by distance to the nearest class prototype accepts roughly
half as many unknowns as maximum-softmax-probability thresholding at matched
true-positive rate --- 0.34 against 0.61 --- and the effect holds across every
variant we tried.

We attributed that to the metric training. We pre-registered the test that
would falsify the attribution: ablate every metric objective on the deployed
embedding, leaving the field-standard configuration, and check whether
prototype rejection still beats softmax.

\textbf{It does, and it does so better.} Removing all metric objectives moves
prototype AUROC from 0.919 $\pm$ 0.033 to 0.949 $\pm$ 0.011 and the false-accept
rate from 0.341 to 0.252. The rejection advantage is a property of the scoring
rule, not of the representation learning we introduced to produce it. On the
same embedding, Mahalanobis distance is statistically indistinguishable from
the prototype rule.

The corrected claim is narrower and more useful than the original: \emph{on
encrypted-traffic embeddings, distance-based rejection substantially outperforms
softmax thresholding for unknown-application detection.} It requires no special
training and applies to anyone's encoder.

\subsection{What did survive of the architecture}

Two things. First, the fused embedding beats its parts: 0.626 $\pm$ 0.001
against 0.616 $\pm$ 0.002 for the sequence embedding alone and 0.613 $\pm$
0.005 for an unlearned concatenation of both. The cross-attention fusion earns
roughly a point of macro-F1, and --- more tellingly --- naive concatenation is
\emph{worse} than the sequence embedding alone, so the learned fusion is doing
something a join does not. Second, the encoder builds a genuinely better task
representation than its input: on CESNET the label is decodable from
$z_{\mathrm{flow}}$ at $+0.741$ above chance against $+0.120$ from raw
features.

Neither is a paper. Both are honest.
